\scanpage{402}\mspage{\textcircled{1}}
% archive sheet number: 401
XV p.~119
\diagram{\begin{tikzcd}[ampersand replacement=\&,column sep=small]
X \arrow[rr] \arrow[dr] \& \& Y \arrow[dl]\\
\& S \&
\end{tikzcd}}
\[
\begin{aligned}
\mathrm{I}^{\circ}/)\quad\text{\textcircled{A}}\quad\mathrm{Loc}(8)&\Longrightarrow\mathrm{Loc}(1\ \mathrm{bis})\\
\text{\textcircled{B}}\quad\mathrm{Loc}(6'\mathrm{a},8,9)&\Longrightarrow
\mathrm{Loc}(1\ \mathrm{ter},\underset{\substack{\vert\\\inserted{\text{\unsure{non|voir} }4\mathrm{b}}}}{4\mathrm{a}},6\mathrm{a})
\end{aligned}
\]
\margin{$\swarrow$ il suffit $=$ $\mathrm{Loc}(6'\mathrm{a},8\deleted{\unsure{)}}\mathrm{bis})$}

Cela implique-t-il \deleted{$\mathrm{Loc}(3\ \mathrm{bis})$} ?
\[
W=W_D,\ \text{\unsure{ou|où} }D=\mathrm{HOT}_{W}\ ?
\]
J'en doute ! (\unsure{Histoire} de triples de Quillen très spéciaux ?)
Et je doute aussi, par suite, que ça implique $\mathrm{Loc}(3\ \mathrm{bis})$.

\figure{\begin{tikzpicture}\draw (0,0) -- (0.75,-1.1);\draw (0.94,-1.34) -- (0.99,-1.43);\end{tikzpicture}}
\[
\deleted{\unsure{\mathrm{II}}}\mathrm{I}\quad
\left\{\begin{aligned}
\mathrm{Loc}(6'\mathrm{a})&\Longrightarrow\mathrm{Loc}(4\mathrm{a}_{\unsure{0|\mathrm{o}}})\\
\mathrm{Loc}(6'\deleted{\unsure{\mathrm{a}}}\mathrm{b})&\Longrightarrow\mathrm{Loc}(4\deleted{\unsure{\mathrm{a}_{0}}}\mathrm{b})\\
\mathrm{Loc}(6\deleted{\unsure{\prime}},8)&\Longrightarrow\mathrm{Loc}(4)
\end{aligned}\right.
\]
\[
\left\{\begin{aligned}
\mathrm{Loc}(7D_{W_{\unsure{\omega|\infty|0}}})&\Longrightarrow\mathrm{Loc}(4\mathrm{b})\\
\deleted{\unsure{\mathrm{Loc}(3}}\ &\\[-2pt]
\deleted{\mathrm{Loc}(7D_{W_{\omega}},} &\\
\mathrm{Loc}(6'\mathrm{a},7,8)&\Longrightarrow\mathrm{Loc}(4)\\
\deleted{\mathrm{Loc}(3\ \mathrm{bis})}&
\end{aligned}\right.
\]
\[
\unsure{\mathrm{II}_{\mathrm{e}}|\mathrm{III}_{\mathrm{e}}}\quad
\left\{\begin{aligned}
\mathrm{Loc}(3\ \mathrm{bis})&\Longrightarrow\deleted{\unsure{\mathrm{D}}}\mathrm{Loc}(4\mathrm{a}_{0})\\
\mathrm{Loc}(3\ \mathrm{bis},8)&\Longrightarrow\deleted{\unsure{\mathrm{D}}}\mathrm{Loc}(4\mathrm{a})\\
\mathrm{Loc}(3\ \mathrm{bis},7,8)&\Longrightarrow\deleted{\unsure{\mathrm{D}}}\mathrm{Loc}(4)
\end{aligned}\right.
\]
\[
\mathrm A\left|\begin{aligned}
\mathrm{Loc}(4\mathrm{a}_{0}):&\ \text{impliqué par }\mathrm{Loc}(6'\mathrm{a}),\ \mathrm{Loc}(3\ \mathrm{bis})\\
\mathrm{Loc}(4\mathrm{a}):&\ \text{impliqué par }\mathrm{Loc}(6'\mathrm{a},8\deleted{\mathrm{bis}})\deleted{\unsure{9}},\ \mathrm{Loc}(3\ \mathrm{bis},8\deleted{\mathrm{bis}})\\
\mathrm{Loc}(4\mathrm{b}):&\ \text{impliqué par }\mathrm{Loc}(6'\mathrm{b})\ (\text{donc par }\mathrm{Loc}(6)),\ \mathrm{Loc}(7D_{W_{\omega}})\\
\mathrm{Loc}(4):&\ \text{impliqué par }\mathrm{Loc}(6\deleted{\unsure{\prime}},8\deleted{\mathrm{bis}}),\
\mathrm{Loc}(6'\mathrm{a},7\deleted{\unsure{D_{W_{\omega}}}},8\deleted{\mathrm{bis}}),\\
&\quad\deleted{[\mathrm{Loc}(3\ \mathrm{bis},6'\mathrm{b},8\mathrm{bis})]}\deleted{\unsure{,}}
\mathrm{Loc}(3\ \mathrm{bis},7\deleted{\unsure{D_{W_{\omega}}}},8\deleted{\mathrm{bis}})
\end{aligned}\right.
\]
